% ===== 原书 p.4 =====
\section{环面 Calabi--Yau 流形及其镜像}

镜像对称原本是对紧 Calabi--Yau 流形考虑的，但后来被推广到了未必紧的 Calabi--Yau 流形、Fano 流形、一般型流形以及奇点等各种对象上。非紧 Calabi--Yau 流形中最重要的例子，当属所谓的环面 Calabi--Yau 流形。之所以说「所谓的」，是因为通常考虑的不是 Ricci 平坦的 K\"ahler 度量而只是 K\"ahler 度量，并且常常考虑的不是环面流形本身，而是其中反典范除子的补集。K\"ahler 流形与其上全纯体积形式之组称为殆 Calabi--Yau 流形（almost Calabi--Yau manifold）。在殆 Calabi--Yau 流形上，存在处处不为零的函数 $f$，使得代替 (4) 而有
\begin{equation}
f^n\frac{\omega^n}{n!}=(-1)^{n(n-1)/2}\left(\frac{\sqrt{-1}}{2}\right)^n\Omega\wedge\overline{\Omega}
\tag{6}
\end{equation}
\noindent 成立。殆 Calabi--Yau 流形上的特殊 Lagrange 子流形，与 Calabi--Yau 流形的情形同样定义。殆 Calabi--Yau 流形上的特殊 Lagrange 子流形关于原来的 Riemann 度量 $g$ 未必极小，但 $\Re\Omega$ 是关于与 $g$ 共形等价的 Riemann 度量 $fg$ 的校准形式，而特殊 Lagrange 子流形是关于这个校准形式的校准子流形。

对 $n$ 维环面 Calabi--Yau 流形 $X$，适当地取直和分解 $M\cong\overline{M}\oplus\Z$ 与格点多面体 $\Delta\subset\overline{M}_{\R}$，便可以把 $X$ 的扇多面体（fan polytope，即扇 $\Sigma$ 的一维锥的原始生成元的凸包）写成 $\Delta\times\{1\}$。关于这组坐标，与 $(0,1)\in N\cong\overline{N}\oplus\Z$ 对应的特征标 $\chi_{(0,1)}:T\to\C^\times$ 唯一地延拓为 $X$ 上的全纯函数 $w_n$。把这个特征的核记作 $T_0:=\Ker\chi_{(0,1)}$，并把与 Lie 代数嵌入 $\iota_*:\overline{M}_{\R}\to M_{\R}$ 对偶的满射记作 $\iota^*:N_{\R}\to\overline{N}_{\R}$，则 $X$ 上 $T_0$ 作用的矩映射，可用 $T$ 作用的矩映射 $\mu:X\to N_{\R}$ 写成 $\mu_0:=\iota^*\circ\mu:X\to\overline{N}_{\R}$。$\mu_0$ 的临界点是 $X$ 的环面不变余维 $2$ 子流形之并 $X^{(n-2)}$，临界值是它经 $\mu_0$ 的像 $\Pi_0:=\mu_0\left(X^{(n-2)}\right)$。此外，把由 $p:=\chi_{e_n}-1=0$ 定义的 $X$ 的除子记作 $\check D$，并定义 $\check Y:=X\setminus\check D$。$p$ 是 $X$ 上的函数，但 $X$ 的典范丛平凡，故可看作反典范丛的截面，于是 $\check Y$ 是 $X$ 的反典范除子的补集。$X$ 由环面不变的仿射开集 $X\subset\C[z_1,\dots,z_n]$ 粘贴而得，具有在各开集上写作 $dz_n\wedge\cdots\wedge dz_n$ 的自然全纯体积形式 $\Omega_X$。
% [Tutor Note: 这里取定的是 $M\cong\overline{M}\oplus\Z$ 与 $N\cong\overline{N}\oplus\Z$，即无上划线的格比带上划线的大一秩；多出的生成元记作 $e_n$（后文 $p=\chi_{e_n}-1$ 用到）。于是 $\iota_*:\overline{M}_{\R}\to M_{\R}$ 是嵌入、$\iota^*:N_{\R}\to\overline{N}_{\R}$ 是满射，$\Delta$ 落在 $\overline{M}_{\R}$ 中，底空间 $B:=\overline{N}_{\R}\times\R^{>0}$ 与 $X$ 同为 $n$ 维]
% [Tutor Note: 原文此处印作「$X\subset\mathbb C[z_1,\dots,z_n]$ を貼り合わせて」，所指应为各仿射图 $X_\sigma\subset\mathbb C^n$（$\sigma$ 取遍扇的极大锥），照录不改]
用它定义 $\check Y$ 的全纯体积形式 $\Omega_{\check Y}:=p^{-1}\Omega_X$。它与 $X$ 上适当的 K\"ahler 形式 $\omega_X$ 在 $\check Y$ 上的限制 $\omega_{\check Y}:=\omega_X|_{\check Y}$ 一起，使 $\check Y$ 成为殆 Calabi--Yau 流形。映射 $\pi_{\check Y}:=(\mu_0,|w|):\check Y\to B:=\overline{N}_{\R}\times\R^{>0}$ 是殆 Calabi--Yau 流形 $\check Y$ 上的特殊 Lagrange 环面纤维化，其判别式集合由 $\Pi_0\times\{1\}$ 给出。
% [Tutor Note: 此处印作 $|w|$，而前文定义的是全纯函数 $w_n$（$\chi_{(0,1)}$ 的延拓）；$w=(w_1,\dots,w_n)$ 与 $w_n$ 的关系原页未说明，疑为 $|w_n|$ 之误或对整组坐标的简写；照录不改]
% [Tutor Note: 「各開集合上で $dz_n\wedge\cdots\wedge dz_n$」中首项在原图亦印作 $dz_n$，同一开集上两次出现同一微分为笔误，通常应作 $dz_1\wedge\cdots\wedge dz_n$；照录不改]

$\check Y$ 的镜像，是具有以二次曲线为纤维的、$\overline{N}_{\C^\times}:=\overline{N}\otimes_{\Z}\C^\times$ 上的纤维化结构的仿射流形
\begin{equation}
Y:=\left\{(\bm{x},u,v)\in\overline{N}_{\C^\times}\times\C^2\mid h(\bm{x})=uv\right\}
\tag{7}
\end{equation}
\noindent 这里，$h(x)$ 是 $\overline{N}_{\C^\times}$ 上的 Laurent 多项式，其 Newton 多面体与 $\Delta\subset\overline{M}_{\R}$ 一致。Laurent 多项式 $h(x)$ 的系数是 $Y$ 的复结构参数，与 $\check Y$ 的辛结构相对应。以下在 $Y$ 上考虑辛几何、在 $\check Y$ 上考虑复几何，因此只要 $Y$ 保持光滑，$h(x)$ 的系数便可任意取。
% [Tutor Note: 原书此句与次句互相矛盾。按数学内容，$h$ 的系数应决定 $\check Y$ 的复结构、并对应 $Y$ 的辛结构（$\check Y$ 正是用 (13) 的粘贴函数依 $h$ 造出的），故疑此句首尾的 $Y$ 与 $\check Y$ 对调；次句「$Y$ 上看辛几何、$\check Y$ 上看复几何」才是作者本意。译文两处均照原样，不擅改]

% ===== 原书 p.5 =====
$\overline{N}_{\C^\times}\times\C^2$ 的自然 K\"ahler 形式
\begin{equation}
\omega_{\overline{N}_{\C^\times}\times\C^2}
:=\sqrt{-1}\partial\overline{\partial}
\left((\log|x_1|)^2+\cdots+(\log|x_{n-1}|)^2+|u|^2+|v|^2\right)
\tag{8}
\end{equation}
\noindent 限制到 $Y$ 上，$Y$ 便成为辛流形；但正如 [AAK] 所指出的，要在 $Y$ 上作出 Lagrange 环面纤维化，把 $Y$ 看作 $\overline{N}_{\C^\times}\times\C$ 沿 $Z\times 0$ 的辛爆破
\begin{equation}
Y:=Bl_{Z\times 0}\left(\overline{N}_{\C^\times}\times\C\right)
:=\left\{(\bm{x},u,[v_0:v_1])\in\overline{N}_{\C^\times}\times\C\times\PP^1\mid h(\bm{x})v_0=uv_1\right\}
\tag{9}
\end{equation}
\noindent 中由 $v_0\neq 0$ 定义的仿射子集，较为方便。这里，$Z:=\{x\in\overline{N}_{\C^\times}\mid h(x)=0\}$。把 $Z\times 0$ 在 $\overline{N}_{\C^\times}\times\C$ 中的管状邻域记作 $U$，把 $\overline{N}_{\C^\times}\times\C$ 上以 $U$ 为支集的、在 $U$ 所含 $Z\times 0$ 的管状邻域上恒等于 $1$ 的光滑函数记作 $\chi$，则对充分小的 $\epsilon$，
\begin{equation}
\omega_\epsilon:=p^*\omega_{\overline{N}_{\C^\times}\times\C}
+\frac{\sqrt{-1}\epsilon}{2\pi}\partial\overline{\partial}
\left(\chi(\bm{x},u)\log\left(|u|^2+|h(\bm{x})|^2\right)\right)
\tag{10}
\end{equation}
\noindent 是 $Y$ 上的辛形式。这里，$p:Y\to\overline{N}_{\C^\times}\times\C$ 是辛爆破的结构射。$Y$ 具有 $S^1$ 作用 $(x,u,v)\mapsto\left(x,e^{\sqrt{-1}\theta}u,e^{-\sqrt{-1}\theta}v\right)$，其矩映射为
\begin{equation}
\mu_Y=\frac{1}{2}|u|^2+\frac{\epsilon}{2\pi}|u|\frac{\partial}{\partial|u|}
\left(\chi(\bm{x},u)\log\left(|h(\bm{x})|^2+|u|^2\right)\right)
\tag{11}
\end{equation}
\begin{equation}
=\begin{cases}
\pi\dfrac{1}{2}|u|^2+\dfrac{\epsilon}{2\pi}\dfrac{|u|^2}{|h_{t,1}(\bm{x})|^2+|u|^2} & \chi\equiv 1\ \text{时},\\[2ex]
\pi\dfrac{1}{2}|u|^2 & \chi\equiv 0\ \text{时}.
\end{cases}
\tag{12}
\end{equation}
% [Tutor Note: (12) 两分支前的因子 $\pi$ 与 (11) 的首项 $\frac12|u|^2$ 不相容，疑为 $\frac12|u|^2$ 之误；$h_{t,1}$ 在前文未定义，应出自 [AAK] 中含参数 $t$ 的族 $h_t$，该参数取定某方向后的记号。两处均照录]
$\mu_Y$ 的临界点由 $S^1$ 作用的固定点 $\{(x,u,v)\in Y\mid h(x)=u=v=0\}$ 给出，临界值为 $\epsilon$。借助自然的投影 $Y\to\overline{N}_{\C^\times}$，$\lambda\neq\epsilon$ 处的辛商 $Y_{\mathrm{red},\lambda}:=\mu_Y^{-1}(\lambda)/S^1$ 与 $\overline{N}_{\C^\times}$ 自然等同。这个等同不保持辛形式，但存在光滑依赖于 $\lambda\in(0,\epsilon)\cup(\epsilon,\infty)$ 的 $\overline{N}_{\C^\times}$ 的微分自同胚族 $\varphi_\lambda$，把约化辛形式 $\omega_{\mathrm{red},\lambda}$ 与 $\overline{N}_{\C^\times}$ 上环面不变的辛形式 $\omega_{\overline{N}_{\C^\times},\lambda}$ 联系起来。此时，把 $\left(\overline{N}_{\C^\times},\omega_{\overline{N}_{\C^\times},\lambda}\right)$ 关于环面作用的矩映射用 $\varphi_\lambda$ 拉回，便给出 $Y_{\mathrm{red},\lambda}$ 的 Lagrange 环面纤维化 $\pi_\lambda:Y_{\mathrm{red},\lambda}\to\overline{M}_{\R}$。这个映射连续延拓到 $\lambda=\epsilon$，但在 $S^1$ 作用的固定点处不再光滑。定义映射 $\pi_Y:Y\to\overline{M}_{\R}\times\R^{>0}$ 为把 $y\in\mu_Y^{-1}(\lambda)$ 送到 $\pi(y)=\left(\pi_\lambda([y]),\lambda\right)$ 的映射，则它给出 $Y$ 上的 Lagrange 环面纤维化，其临界值由 $\Pi_\epsilon\times\{\epsilon\}$ 给出。这里，$\Pi_\epsilon:=\pi_\epsilon(Z)$ 是经 $\varphi_\epsilon$ 修正后的 $Z$ 的 amoeba。

与限制 (8) 所得的辛形式相比，(10) 定义的辛形式 $\omega_\epsilon$ 的优点之一，是在 $Z\times 0$ 的管状邻域 $U$ 之外与 $\overline{N}_{\C^\times}\times\C$ 的标准辛形式一致。因此，对与 $Z$（更确切地说，与 $U$ 到 $\overline{N}_{\C^\times}$ 的投影的像）不相交的 $\overline{N}_{\C^\times}$ 的 Lagrange 子流形 $\overline{L}$，$L:=\overline{L}\times\R^{>0}\subset\overline{N}_{\C^\times}\times\C$ 自然地确定 $Y$ 的 Lagrange 子流形。

% CHUNK c02 DONE pages p4–p5
